% Ce fichier n'est jamais servi aux utilisateurs : il sert uniquement, au
% moment de la construction de l'image Docker, à faire télécharger et mettre
% en cache à Tectonic TOUS les packages réellement utilisés par les deux
% templates de l'application (genererTex et genererTexQcm), pour que les
% conteneurs qui tournent ensuite en production n'aient plus jamais besoin
% du réseau pour compiler une interro ou une feuille de QCM.
\documentclass[12pt]{article}
\usepackage[utf8]{inputenc}
\usepackage[T1]{fontenc}
\usepackage[french]{babel}
\usepackage{amsmath,amssymb}
\usepackage{tcolorbox}
\usepackage{fancyhdr}
\usepackage{forloop}
\usepackage{enumitem}
\usepackage{multicol}
\usepackage[margin=2cm]{geometry}
\pagestyle{fancy}
\fancyhf{}
\lhead{Préchauffage}
\chead{Cache}
\rhead{Tectonic}
\newcounter{qnum}
\newcounter{linectr}
\newcommand{\reponse}[1]{
  \setcounter{linectr}{0}
  \forloop{linectr}{0}{\value{linectr} < #1}{\par\vspace{4mm}\hrulefill}
}
\begin{document}

\noindent\textbf{Nom :}\underline{\hspace{3.5cm}} \hfill \textbf{Prénom :}\underline{\hspace{3.5cm}} \hfill \textbf{Classe :}\underline{\hspace{2cm}}

\stepcounter{qnum}
\noindent\textbf{Question \theqnum.} Exemple $x^2 + 1 = 0$.
\begin{tcolorbox}[colback=gray!10]
Corrigé d'exemple.
\end{tcolorbox}
\reponse{3}

\small
\begin{multicols}{2}
\begin{enumerate}[leftmargin=2.4em, itemsep=4mm, label=$\square$~\alph*)]
\item[$\blacksquare$~a)] \textbf{Bonne réponse}
\item Mauvaise réponse
\columnbreak
\item Mauvaise réponse
\item Mauvaise réponse
\end{enumerate}
\end{multicols}

\end{document}
